% Generated draft: requires mathematical and source review.
\documentclass[11pt]{article}
\usepackage[margin=25mm]{geometry}
\usepackage{fontspec}
\setmainfont{Latin Modern Roman}
\usepackage{amsmath,amssymb,mathtools,unicode-math}
\setmathfont{Latin Modern Math}
\usepackage{enumitem,needspace,xurl,hyperref}
\setlength{\parindent}{0pt}
\setlength{\parskip}{5pt}
\setlength{\emergencystretch}{3em}
\newcommand{\N}{\mathbb{N}}
\newcommand{\Z}{\mathbb{Z}}
\newcommand{\Q}{\mathbb{Q}}
\newcommand{\R}{\mathbb{R}}
\newcommand{\C}{\mathbb{C}}
\newcommand{\F}{\mathbb{F}}
\newcommand{\A}{\mathbb{A}}
\newcommand{\PP}{\mathbb{P}}
\newcommand{\E}{\mathbb{E}}
\newcommand{\eps}{\varepsilon}
\newcommand{\dd}{\,\mathrm{d}}
\begin{document}
\section*{UnsolvedMath ID: 30002571}
\textbf{Source problem code:} catalogue entry
\subsection{Definitions and mathematical statement}
\textbf{Research question.} The supplied record asks: ``Do the common structures of complex living systems provide a paradigm for deriving life-science models, and can applications reveal hallmarks of a mathematical theory of living systems?''

\textbf{Unresolved definitions and scope.} The record does not define a class of complex living systems, the structures they are required to share, a mathematical representation of those structures, or a procedure for deriving models. It also does not specify the relevant applications, the meaning of a hallmark, or criteria for determining whether a paradigm or mathematical theory has been obtained. No mathematical domains, hypotheses, quantifiers, or boundary conditions are supplied. Consequently, this is a research-programme question; a precise quantified mathematical assertion cannot be extracted without additional specifications. The wording above is retained without introducing such specifications.
\subsection{Short English statement}
The question asks whether shared structures of living systems can guide the derivation of life-science models, and whether applications can identify characteristic features of a mathematical theory of living systems. The supplied record leaves the central concepts and success criteria undefined, so it does not determine a precise mathematical proposition.
\subsection{Sources}
\begin{itemize}
\item Supplied local registry record: permanent UnsolvedMath ID \(30002571\), title ``Mathematical Principles for Modelling Living Systems,'' collection ``Oberwolfach Reports Open Problems,'' source code \texttt{catalogue entry}. The original, clean, and main statements agree. The registry labels the statement \texttt{exact} and describes it as retained byte-for-byte, clear, self-contained, and publication-ready; these are supplied editorial qualifications, not independent mathematical verification. No individual proposer is identified.
\item Supplied report citation: \emph{Mini-Workshop: Mathematical Models for Cancer Cell Migration (2015)}, Oberwolfach Reports, DOI 10.4171/owr/2014/19, pp.~1075--1109. \url{https://doi.org/10.4171/owr/2014/19}. The background describes Oberwolfach Reports as documenting mathematical workshops and their open-problem discussions. The registry gives proposed year 2014, whereas the citation labels the workshop 2015; this discrepancy is unresolved.
\item Supplied literature source, labeled ``Theory of living systems review'': \url{https://doi.org/10.1016/j.plrev.2010.12.004}.
\item Supplied literature source, labeled ``Open problems in mathematical biology'': \url{https://arxiv.org/abs/2206.09516}.
\item Literature assessment dated exactly \texttt{2026-08-21}: the supplied curation labels the problem open and states that mathematical biology has produced useful frameworks and proposed hallmarks, but that a generally accepted paradigm or unified mathematical theory of living systems has not been established. This is supplied dated curation derived from the OWR material, not an independently verified assessment. The source explicitly requests independent verification of time-sensitive status and cited claims.
\end{itemize}
\end{document}
